\section{Future work}
\label{sec:future}

\subsection{Identifiability first}
Before individual heads can be interpreted or used in audits, we need to
ensure that they reproduce across training runs, or else identify the
subset of heads that do. In our current models only the first layer
shows partial reproducibility, with a single topic-and-genre head
recurring across SONAR runs; understanding why that head recurs and how
to extend reproducibility to other heads is our primary focus. Classical
identifiability theory points to two approaches:
\begin{enumerate}
	\item selecting reproducible heads or runs via unsupervised
	cross-seed agreement~\cite{udr2020} (in latent-variable models, three
	or more conditionally independent views of a discrete variable
	identify its distribution up to permutation when the
	class-conditional distributions satisfy a rank
	condition~\cite{allman2009}, with individual heads serving as
	candidate views here); and
	\item breaking the symmetry of the marginal distribution~\cite{locatello2019}
	using paired observations such as adjacent sentences in a
	document~\cite{pcl2017,slowvae2021,locatello2020,ivae2020,zimmermann2021},
	since standard sparse-coding uniqueness guarantees require lossless
	reconstruction~\cite{hillar2015}.
\end{enumerate}

Initial experiments on the 44.1M-segment sentence-split SONAR cache
(\texttt{mc4\_sent16}; working notes) show that a sentence embedding
$\mathbf{x}_t$ linearly explains $10.4\%$ of the variance of the
following sentence $\mathbf{x}_{t+1}$
($\mathrm{FVU}_{\mathrm{next}}{=}0.896$), and that adjacent sentences are
strongly aligned in direction (raw cosine AUC $0.923$ vs.\ $0.685$ in
whitened distance), yet heads from single-sentence models switch entries
on $92.2$--$98.1\%$ of adjacent pairs. Directly penalizing entry
switches across adjacent sentences on 38 designated slow heads collapses
their usage entropy to $0.04$--$0.09$ bits for penalty weights
$s_{\mathrm{pair}} \ge 10^{-4}$. Replacing that penalty with a Fano
lower bound on each slow head's mutual information across the sentence
pair prevents collapse and produces a block of slow heads that
reproduces across three seeds, though not head by head
(Section~\ref{sec:prelim-slow}). A first attempt to tell the slow heads
apart, an ordered penalty on dependence between their partitions,
identifies one head at full weight and strips the rest of their slow
content. Immediate next steps include: separating the slow heads, which
agree with one another nearly as well as with another seed's, by a
criterion that keeps what the block carries jointly; extending
the objective to sentence triples to test the three-view
conditional-independence rank conditions on held-out
data~\cite{allman2009}; characterizing what the slow heads encode beyond
language; and applying the paired-view
objective to a causal language model over multi-turn conversations (such
as Gemma~3 1B instruct). In parallel, replacing the shared decoder in
the 40M-token GPT-2 concatenated-head model with per-layer decoders
reduces cross-layer same-head coupling from $97\%$ to $6\%$ while
improving $\mathrm{FVU}$ from $0.186$ to $0.169$ (pending a
width-matched shared-decoder control).

\subsection{Steering at scale}
We have compared model-derived steering directions against supervised
vectors on the models' own partitions (working notes) and, for the
softmax model, on language steering (Section~\ref{sec:prelim-steer}).
Several steering evaluations remain open: running the language-steering
comparison on hard codes; evaluating AxBench-style steering toward
external concepts against supervised vectors and SAE latents at tuned
intervention strengths (initial Gemma~3 1B and Gemma~4 A4B models with
per-layer decoders are currently training); testing bilinear
eigenfeatures as steering directions; measuring whether an entry swap
produces consistent semantic shifts across contexts; and trading off the
single-layer model's $1.79\times$ higher reconstruction error against its
near-perfect target realization and lower collateral.

Reducing input-space collateral requires widening classifier margins
rather than merely finding better steering directions: in the hard stack,
heads sit near their decision boundaries, and even oracle supervised
vectors reduce collateral flips by at most 0.1 relative to a random step
of the same norm. Reducing the number of entries per head increases those
margins. In a five-layer stack of binary ($k{=}2$) heads, a step of norm
$0.25\lVert\mathbf{x}\rVert$ is predicted to flip 13\% of same-layer
heads to first order (vs.\ 30\% for $k{=}32$), and empirically flips
17--19\% of other heads across all directions (vs.\ 66\% for the
$k{=}32$ stack), matching that lower collateral at the 32-entry stack's
step count under a matched per-head $\mathrm{KL}_m$ weight. Because that
binary run used a single seed and half the batch size, however, a
controlled comparison isolating $k$ is still needed, with the
single-layer linear model \texttt{g160top1} (Section~\ref{sec:prelim-depth})
serving as the baseline for low-collateral addressing.

\subsection{Task reconstruction}
For heads that reproduce across seeds, we aim to identify the classification
task each head performs by measuring mutual information with metadata
labels (such as topic and genre), generating contrastive natural-language
descriptions that distinguish an entry from its sibling entries (scored by
blinded $k$-way assignment accuracy against a $1/k$ chance baseline), and
using parallel translations to separate content-sensitive heads from
surface-form heads.

\subsection{Disjoint heads}
Because summed models fail to partition dictionary coordinates across
heads and penalizing $\omega$ on dictionary weights merely pushes atom
energy into the decoder's near-null space (leaving decoded overlap
8--10\% above the coordinate-shuffle null; working notes;
cf.~\cite[§5.1]{locatello2019}), any learned overlap penalty must operate
on decoded or direct atoms where the decoder cannot absorb it.
Alternatively, concatenated heads enforce disjoint dictionary blocks by
construction and already outperform direct atoms on GPT-2 when each
head's coordinate slice is wider than $k$ (Section~\ref{sec:prelim-gpt2}).

\subsection{Transcoder adaptation}
To adapt the architecture into a transcoder that predicts a sublayer's
output $Y$ from its input $X$~\cite{dunefsky2024}, we tested two pilot
designs on GPT-2 block~8's MLP (10M tokens, concatenated heads; working
notes). A dual-decoder stack that reconstructs both $X$ and $Y$ from the
shared residual accumulator beats the linear ridge baseline on $Y$
($\mathrm{FVU}_Y$ $0.769$, $28.3\%$ MLP-splice loss recovered; ridge on
SAE latents reaches $0.612$) only when the loss weight on $Y$ is raised
to $\lambda \ge 3$ ($\mathrm{FVU}_Y$ $0.763 \to 0.729$, $28.6\% \to
32.1\%$ loss recovered), which severely degrades input reconstruction
($\mathrm{FVU}_X$ $0.284 \to 0.475$). Conversely, a target-free
forwarding mode (\texttt{forward="cond"}) that does not require $Y$ at
inference time reaches only $\mathrm{FVU}_Y = 0.855$ ($10.0\%$ loss
recovered), with deeper layers increasing error ($0.823 \to 0.855$). A
viable multi-layer transcoder will require a residual recurrence that
iteratively refines predictions of $Y$ without access to the ground-truth
target during the forward pass.
